\chapter{Introduction}
\label{Chapter1}

\section{Motivation and Problem Statement}

\textit{[TODO: Describe the context of Android malware and why detection is
important. Explain the problem this thesis addresses -- evaluating how well
Large Language Models (LLMs) perform at Android malware detection -- and why
existing evaluations are insufficient or incomplete (research gap identified
in the SLR).]}

\section{Research Questions and Objectives}

\textit{[TODO: State the specific research questions the thesis answers, and
the objectives of the evaluation framework, e.g. comparing LLM-based methods
against baselines such as the LAMD framework under consistent datasets and
metrics.]}

\section{Contributions}

\textit{[TODO: Summarize the contributions of this thesis, e.g. (1) a
systematic literature review of LLM-based Android malware detection
approaches, (2) an evaluation framework/methodology for comparing candidate
methods, (3) an empirical comparison of selected methods including one built
on top of the LAMD framework.]}

\section{Thesis Structure}

The remainder of this thesis is organized as follows.
Chapter~\ref{Chapter2} reviews background concepts on Android malware
detection and LLM-based approaches.
Chapter~\ref{Chapter3} presents the first phase of the thesis, a systematic
literature review of LLM-based Android malware detection, and identifies the
research gaps that motivate the evaluation.
Chapter~\ref{Chapter4} presents the proposed evaluation methodology.
Chapter~\ref{Chapter5} reports and discusses the experimental results.
Chapter~\ref{Chapter6} concludes the thesis and outlines future work.
